\section{Business Implications, Forecast Output, and Limitations}

\paragraph{From forecast to business action.}
The model is useful because it translates historical patterns into forward-looking business risks. The main forecast result is not only a Kaggle score; it is a structured view of where the business should pay attention: conversion quality, seasonal demand planning, H2 margin volatility, Streetwear inventory exposure, wrong-size returns, and COD cancellation risk.

\begin{table}[h]
\centering
\caption{Business findings translated into recommended actions and monitoring metrics.}
\label{tab:business_actions}
\begin{tabular}{p{0.25\linewidth}p{0.32\linewidth}p{0.28\linewidth}}
\toprule
Finding & Recommended action & Metrics to monitor \\
\midrule
Conversion declined despite higher sessions & Audit landing pages, product pages, cart flow, and checkout completion before increasing traffic spend & Conversion proxy, checkout completion, revenue per session \\
Organic search has high volume but weak commercial efficiency & Improve keyword-intent matching and landing-page relevance; retarget high-intent browsers & Organic revenue per session, organic conversion proxy \\
Revenue peaks in April--June & Prepare inventory, marketing, and logistics capacity before seasonal peaks & Monthly demand forecast, stockout days, fulfillment SLA \\
Streetwear dominates revenue and stockout exposure & Prioritize replenishment for high-demand Streetwear SKUs, sizes, and colors & Fill rate, days of supply, sell-through rate \\
COGS ratio is volatile in H2 and promotion-heavy months & Plan margin buffers and evaluate promotions by gross margin, not only revenue uplift & COGS/Revenue ratio, promo margin, discount rate \\
Wrong-size returns create refund leakage & Improve size guides, fit recommendations, and product descriptions & Wrong-size return count, refund amount, return rate \\
COD has higher cancellation risk & Confirm COD orders, limit COD on high-risk baskets, and incentivize prepaid payment & COD cancellation rate, prepaid share, cancelled value \\
\bottomrule
\end{tabular}
\end{table}

\paragraph{Forecasting output.}
The final forecast achieved a Kaggle MAE of \metric{667,139.05}. More importantly, the forecast is built around two business equations:
\[
  \widehat{\text{Revenue}} =
  \widehat{\text{Demand level}}
  \times
  \widehat{\text{Seasonal allocation}},
  \qquad
  \widehat{\text{COGS}} =
  \widehat{\text{Revenue}}
  \times
  \widehat{\text{Margin ratio}}.
\]
This framing is important because a revenue-only forecast can hide margin compression. A promotion-heavy day may look healthy in Revenue while still carrying high COGS risk.

\paragraph{Validation and sanity checks.}
Validation is time-aware. The workflow uses long-horizon and year-based folds to mimic forecasting more than one year ahead, rather than random folds that would leak adjacent temporal structure. Candidate forecasts are checked for row count, date range, non-negative values, period totals, maximum daily values, and COGS/Revenue ratio ranges before submission.

\paragraph{Limitations.}
Future operational facts such as actual 2023--2024 orders, returns, inventory, reviews, and traffic are unavailable at forecast time. The model therefore uses historical priors for these drivers. This makes the pipeline reproducible and avoids unavailable future inputs, but it limits the ability to react to new campaigns, supply shocks, or acquisition changes after 2022.

\paragraph{Summary.}
The recommended strategy is not to increase traffic or discounts indiscriminately. The business should first convert existing traffic more effectively, protect the core Streetwear category before seasonal peaks, and manage margin risk during promotion-heavy H2 periods. The forecast supports this strategy by separating demand seasonality from COGS ratio behavior.
